\chapter*{Abstract}
\addcontentsline{toc}{chapter}{Abstract}

Large Language Models (LLMs) have demonstrated remarkable capabilities in natural language processing tasks. However, they are prone to generating fluent but factually incorrect or logically inconsistent responses, a phenomenon known as ``hallucination.'' Detecting hallucinations without access to ground-truth answers remains a fundamental challenge, as the most informative detection signals---such as comparing outputs against known correct answers---are inherently unavailable during real-world inference.

This thesis presents LOGIC-HALT (LOGical Inconsistency and Compression-based HALlucination Testing), a multi-signal hallucination detection framework that operates entirely without ground-truth answers at inference time. The system combines five complementary signals through a Bayesian-optimized fusion layer: (1) self-verification, where the model is asked to fact-check its own answer via Natural Language Inference (NLI), (2) cross-response inconsistency measured through pairwise NLI contradiction analysis using DeBERTa-v3-large (435M parameters), (3) token entropy computed from log-probabilities, (4) pairwise Normalized Compression Distance (NCD) between response variants, and (5) answer-level minority penalty from majority voting across extracted answers. These signals are generated by querying the target LLM with semantically equivalent question variants produced through metamorphic transformations---including paraphrasing, logical negation, variable substitution, premise reordering, and redundant context injection---each validated for semantic preservation via Sentence-BERT cosine similarity.

A critical contribution of this work is the identification and correction of a ground truth leakage problem in the original evaluation framework. The initial system incorporated a Ground Truth Contradiction signal that compared model outputs against known correct answers, accounting for 63.8\% of the fusion weight and producing artificially inflated metrics (Precision = 0.977). Since this signal is unavailable at inference time, we replaced it with inference-time proxy signals and re-optimized all hyperparameters using Bayesian optimization (Optuna TPE sampler, 500 trials, 13 hyperparameters).

Evaluation on the TruthfulQA benchmark (817 questions) yields an F1 score of 0.727 (Precision = 0.602, Recall = 0.917) under precision-balanced optimization in fully reference-free conditions. When optimized purely for F1 without precision constraints, the system achieves F1 = 0.735 with perfect recall but reduced precision (0.581). A systematic ablation study across 8 signal combinations reveals that answer-level majority voting is the strongest individual signal (39.19\% optimized weight), and that Negation Contradiction is the only additional signal that improves precision (0.578 $\rightarrow$ 0.606). The system is implemented with GPU optimization for NVIDIA RTX 5070 (12\,GB VRAM, CUDA 12.8) with FP16 inference and batch NLI processing, and includes a Flask-based web interface for real-time analysis.

The findings demonstrate that while multi-signal consistency analysis achieves high recall for hallucination screening, the fundamental limitation of consistency-based approaches---that consistency does not imply correctness---caps precision at approximately 60\% without external knowledge sources. This establishes a clear direction for future work integrating retrieval-augmented verification and cross-model consistency signals.

\vfill
\textbf{Keywords:} Large Language Models, Hallucination Detection, Metamorphic Testing, Natural Language Inference, Information Theory, Bayesian Optimization, TruthfulQA, Normalized Compression Distance.
\clearpage
